\documentclass{article}
\begin{document}

\date{}

\noindent Dear Editor,

\vspace{1em}

We submit the enclosed manuscript entitled ``A VOF-Guided Neural Surrogate for Position-Conditioned Flow Prediction in Argon-Stirred Ladle Refining'' for consideration in \textit{Engineering Applications of Computational Fluid Mechanics}.

This manuscript is a substantially revised and restructured version of a previous submission to \textit{Engineering Applications of Computational Fluid Mechanics} (prior manuscript ID available upon request). We are grateful for the reviewers' detailed feedback, which helped us identify major issues in the original train/test partitioning, the interpretation of physics-informed training, and the baseline comparison. The present manuscript has therefore been fully reformulated in terms of methodology, evaluation protocol, and narrative structure.

The major changes are summarized as follows:

\begin{enumerate}
\item \textbf{VOF-guided domain restriction.} The original manuscript attempted full-domain PINN-style training without explicitly addressing the multiphase density range. The revised manuscript restricts learning to a VOF-filtered steel-containing computational subdomain ($\alpha_{\rm steel}>0.01$), where the retained bulk flow is dominated by the steel phase. This physically motivated restriction improves learning stability and substantially improves prediction performance compared with unfiltered training.

\item \textbf{Position-conditioned prediction with explicit train/test protocols.} A binary injection-position parameter $m\in\{0,1\}$ is introduced to enable a single surrogate model for both center-bottom and eccentric-bottom blowing. Three evaluation protocols are now clearly defined: center-only prediction, position-conditioned prediction, and unseen-position testing. The revised train/test design removes the ambiguity present in the previous submission.

\item \textbf{POD-ROM benchmarking.} Proper orthogonal decomposition is included as a strong reduced-order-model baseline. The results show that POD-ROM achieves high accuracy for center-blowing cases but does not directly transfer to eccentric-bottom blowing in its center-trained form. This comparison clarifies the trade-off between single-configuration reduced-order modeling and the proposed position-conditioned neural surrogate.

\item \textbf{Expanded validation and reproducibility.} The revised manuscript reports per-condition and per-component validation metrics, including $R^{2}$ and RMSE. The surrogate modeling code, trained model weights, result files, paper figures, and a representative CFD export sample are made available for reproducibility.
\end{enumerate}

A companion manuscript by the authors focuses on the CFD-DPM bubble--inclusion contact-screening algorithm and inclusion-transport analysis. The present manuscript is distinct in scope: it develops and validates a VOF-guided neural surrogate for rapid position-conditioned liquid-steel flow-field prediction. All CFD settings required to understand and assess the surrogate dataset are summarized in the present manuscript, and the companion work is cited only for methodological continuity.

We believe that the revised manuscript represents a meaningful contribution to CFD-informed surrogate modeling for industrial ladle metallurgy. It demonstrates that VOF-guided domain restriction combined with position-conditioned neural networks can provide a fast and practical surrogate for three-dimensional flow-field prediction under different argon-stirring configurations.

We confirm that this manuscript has not been published elsewhere and is not under consideration by another journal. All authors have approved the manuscript and agree with its submission to \textit{Engineering Applications of Computational Fluid Mechanics}.

\vspace{1em}
\noindent Sincerely,\\
The Authors

\end{document}
